\section{Training, validation, and complete work}\label{sec:training69}
The fixed-feature proposition is now executed as its own numerical arm, rather than being attributed to the earlier hidden-layer fits. For $x\in[0,1]^d$, compare the two nonnegative feature maps
\[
 \phi_{jh}^{\rm R}(x)=d^{-1}(1-|4x_h-j|)_+,
 \qquad
 \phi_{jh}^{\rm B}(x)=d^{-1}{4\choose j}x_h^j(1-x_h)^{4-j},
 \quad j=0,\ldots,4.
\]
Both sum to one and have $5d$ coefficients. The first is an affine--ReLU hat construction; its circuit and its hat basis are one method. The second is the conventional degree-four Bernstein construction. Both readouts receive the same purchased states, boundary Bellman enclosures and approximate labels, the same certificate-aware objective, box constraint, penalty, iteration cap and weight lattice. A label-loss ReLU control retains the same features but changes the objective. The quadratic control receives the same information, with its different coefficient count reported. Prediction-free controls incur no unused training cost.

After quantizing the returned weights, we recompute $F(w)$ and the box dual gap $D(w)$ exactly from rational contexts and features. This certifies the achieved iterate even if the finite optimizer reports nonconvergence. The rounding allowance is checked against actual rounded training proposals. Future policy returns still require the outward screen; neither a small training loss nor the optimizer's status can authorize an action.

Validation uses a disjoint, predictor-independent stream of uniform acquired-state contexts. A failed first screen invokes a declared complete-action mesh with a finite root-query cap $K(z)$. This recovery and its descendants are executed. Its context law is not the distribution of states visited by recursive deployment, so the distributional validation conclusion is not silently transferred to those states.

\begin{proposition}[A finite catalogue of complete work]\label{prop:complete-work69}
Condition on a frozen finite family of predictors and a predictor-independent validation law. Suppose a failed screen invokes a deterministic recovery whose root-query count is at most $K(z)$, and every root query costs at most $W(z)$ under the declared descendant and arithmetic contract. If $0\le W(z)K(z)\le C_{\max}$, then with probability at least $1-\alpha$, simultaneously for $M$ fixed predictors,
\[
 \mathbb E[\text{recovery work}_\theta]
 \le\frac1n\sum_{i=1}^n W(z_i)K(z_i)\ell_\tau(z_i,\theta)
       +C_{\max}\sqrt{\frac{\log(M/\alpha)}{2n}}.
\]
Training, context acquisition, prediction, loading, identity checks and output costs are additional terms in complete work. Replacing $W$ by one gives a root-query result, not a bound on elapsed process time.
\end{proposition}
\begin{proof}
On a failed screen the clipped loss equals one and recovery work is at most $WK$; on success recovery work is zero. Apply the bounded-sum inequality to $WK\ell_\tau\in[0,C_{\max}]$ and take a union bound. The separately named service costs are not recovery work.
\end{proof}

The executed validation uses the root-query version. Complete process clocks remain separate measurements. Each reuse process fits and serializes once, then executes sixty-four declared blocks. Model loading checks the source identity, task, horizon, acquisition precision and stored hash. Block sixteen deliberately presents a stale acquisition identifier and uses conventional verification. An early persistent crossing requires cumulative measured work below the matched adaptive process by block forty-eight and at every later block through sixty-three, using zero-based indices. A last-block crossing does not qualify. Four process repetitions describe timing variability; they do not establish an infinite-horizon crossing or a hardware-invariant probability.

For two controls, a trained two-output proposal supplies a paid feasible upper witness while the original triangular lower cover remains intact. Projection and downward rounding precede its original-law evaluation.
\begin{proposition}[Optional vector witnesses]\label{prop:vector-witness69}
Suppose a triangular action cover has a valid global lower bound $L$ and feasible evaluated witnesses with upper bound $U$. Adding any finite number of feasible, correctly evaluated witnesses and replacing $U$ by the minimum of all their upper bounds preserves containment of the original feasible infimum. Every returned witness with $U-L\le\delta$ has original Bellman regret at most $\delta$. The conclusion is independent of the training method and of whether an added witness is a vertex of the lower cover.
\end{proposition}
\begin{proof}
The unchanged cover retains its lower bound. Every evaluated feasible action supplies an upper bound for the infimum; their minimum remains an upper bound. The returned witness has cost at most its own upper endpoint.
\end{proof}

The vector experiment compares adaptive critical-triangle and uniform largest-edge refinement with two-output ReLU and quadratic proposals under the same certificate. Uniform refinement changes the triangle selected for a split, not the global lower-bound calculation. Four-date neural proposals transfer three-date fitted weights. Vector fitting uses action-label loss and is not relabeled as the scalar certificate-aware training theorem. The complete development companion gives the full training, validation, reuse and witness specification.
